\documentclass[11pt]{article}
\usepackage[T1]{fontenc}
\usepackage{lmodern}
\usepackage[margin=1.05in]{geometry}
\usepackage{amsmath,amssymb,amsthm}
\usepackage{booktabs,array}
\usepackage{microtype}
\microtypesetup{expansion=false}
\usepackage{xcolor}
\definecolor{linknavy}{RGB}{25,55,125}
\usepackage[colorlinks=true,linkcolor=linknavy,citecolor=linknavy,urlcolor=linknavy]{hyperref}
\usepackage{orcidlink}
\usepackage{fancyhdr}
\hypersetup{
  pdftitle={A quadratic family of integral two-torsion classes in a Huneke-Wiegand Koszul presentation},
  pdfauthor={Felipe Santibanez-Leal},
  pdfsubject={Explicit integral two-torsion, independent triangle classes, and a direct summand of a Koszul cokernel},
  pdfkeywords={Koszul differential, integral cokernel, two-torsion, direct summand, numerical semigroups, exact certificates}
}
\newtheorem{theorem}{Theorem}[section]
\newtheorem{proposition}[theorem]{Proposition}
\newtheorem{lemma}[theorem]{Lemma}
\newtheorem{corollary}[theorem]{Corollary}
\theoremstyle{definition}
\newtheorem{definition}[theorem]{Definition}
\theoremstyle{remark}
\newtheorem{remark}[theorem]{Remark}
\newcommand{\Z}{\mathbb Z}
\newcommand{\coker}{\operatorname{coker}}
\newcommand{\supp}{\operatorname{supp}}
\newcommand{\src}[3]{#1[#2;#3]}
\newcommand{\VersionDOI}{10.5281/zenodo.22342976}
\newcommand{\ConceptDOI}{10.5281/zenodo.22342975}
\newcommand{\Exp}[2]{\href{https://github.com/fsantibanezleal/CAOS_RESEARCH/tree/main/problems/commutative-algebra/huneke-wiegand/experiments/#1}{#2}}
\newcommand{\Repo}{\href{https://github.com/fsantibanezleal/CAOS_RESEARCH}{CAOS\_RESEARCH}}
\setlength{\headheight}{14pt}
\pagestyle{fancy}
\fancyhf{}
\fancyhead[L]{\small CAOS Research Preprint}
\fancyhead[R]{\small Quadratic Integral Two-Torsion \textperiodcentered\ v0.01}
\fancyfoot[C]{\small\thepage}
\renewcommand{\headrulewidth}{0.4pt}
\fancypagestyle{plain}{\fancyhf{}\fancyfoot[C]{\small\thepage}\renewcommand{\headrulewidth}{0pt}}

\title{A quadratic family of integral two-torsion classes\\
in a Huneke--Wiegand Koszul presentation}
\author{Felipe Santiba\~nez-Leal\,\orcidlink{0000-0002-0150-3246}\\
\small CAOS Research programme, \texttt{github.com/fsantibanezleal/CAOS\_RESEARCH}}
\date{2026-09-05\\[2pt]\normalsize Version 0.01}

\begin{document}
\maketitle
\thispagestyle{plain}
\begin{center}\small
\fbox{\parbox{0.92\textwidth}{\centering
\textbf{PREPRINT} \textperiodcentered\ not peer reviewed \textperiodcentered\ CAOS Research programme\\
Homological companion in the Huneke--Wiegand extension record\\
Version 0.01 \textperiodcentered\ 2026-09-05 \textperiodcentered\ licensed CC BY 4.0\\
Version DOI: \href{https://doi.org/\VersionDOI}{\VersionDOI}\\
Concept DOI: \href{https://doi.org/\ConceptDOI}{\ConceptDOI} (resolves to latest)\\
Sources and exact certificates: \Repo\\[2pt]
\textbf{Scope.} A quadratically growing direct summand of integral
two-torsion, not a description of the whole connecting quotient.
}}
\end{center}

\begin{abstract}
We study an explicit integral Koszul presentation arising in the graded-syzygy
analysis of a previously constructed Huneke--Wiegand family. For every integer
$p\geq8$, its cokernel contains a direct summand isomorphic to
$(\Z/2)^{\lfloor((p-2)^2+3)/12\rfloor}$. The classes are indexed by all
strictly increasing nonnegative triples of sum $p-2$. Reflection-symmetric
interval potentials give explicit original integral sources for twice each
class. Relative parity functionals, supported on ten to twelve original rows,
detect the classes independently and vanish on every complete connecting
image; extension over the remaining target gives a nonconstructive retraction.
A previously tracked four-row vector is integrally congruent to one selected
class, with an explicit source for the difference. We also give its full
twice-source, two short kernel-domain corrections, and a neighboring integral
potential basis of rank $\binom{p-1}{2}$ with at most seven-row images.
Independent finite campaigns check the signed identities and complete-sector
parity certificates, preserving all tested sources in lossless archives.
A failed local search is retained as an exact locality obstruction.
The uniform theorem follows from signed reconstruction, exhaustive source
classification, and an elementary counting recurrence, not finite rank
extrapolation. The full cokernel, an upper bound, and identification with
earlier isolated subpresentations remain unresolved.
\end{abstract}

\section{Introduction and scope}

Exact linear algebra can certify a finite presentation without revealing why
the same phenomenon should persist with a parameter. This distinction matters
for the connecting maps studied here: a displayed target vector, an integral
multiple in the image, and a nonzero quotient class are different assertions.
The source identity and relative parity argument close the second and third
of these obligations for a quadratically growing independent family of target
classes. Their subgroup is a direct summand, although its complement remains
undetermined.

The underlying family and its conductor-fiber-cone and graded-syzygy context
were developed in the first paper of this programme \cite{CAOSMain}. We make
the present matrix problem self-contained by defining its complete integral
presentation below. The new result is therefore checkable without reproducing
the semigroup construction. Its application to that family uses the
presentation identified in the earlier record; it does not repeat or claim
priority for the discovery of a counterexample addressed by that programme.

Signed Koszul strands admit useful combinatorial descriptions. For semigroup
quotients, Bruns and Herzog provide a relative simplicial viewpoint
\cite{BrunsHerzog}. Here the arguments use explicit oriented bases and
low-exterior complements. We do not infer a simplicial or projective-plane
model from a matrix with entries $0,1,-1$, and no unidentified topological
model is a premise of the proof. Potentials and signed complements are
classical tools; the contribution is the explicit triangle family in this
presentation, including its full sources, all boundary corrections, and
complete relative functionals. The general fact that vector-space functionals
extend is used only to deduce splitting; it is not claimed as a new method.

Statements proved deductively are marked [D]; finite exact checks are marked
[MV]. Remaining research questions are marked [C]. The annihilator theorem
alone says \emph{order dividing two}. The additional nonvanishing proof in
Section~\ref{sec:parity} upgrades the tracked class to order exactly two.
Section~\ref{sec:family} constructs all triangle classes, proves their exact
relation lattice and count, and derives the direct-summand corollary.
These are lower-bound and partial-decomposition results, not the whole quotient.

\paragraph{Independent validation boundary.}
The generic parity, independence and splitting arguments below are deductive
and have undergone independent paper review. The separate full-sector audit
for all 27 triangles at five declared parameters also passes.
Section~\ref{sec:validation} distinguishes these finite certificates from
the exhaustive all-parameter proof; no finite rank pattern is extrapolated.

\section{The complete integral presentation}\label{sec:matrix}

All intervals in this paper contain integers, and an empty interval contributes
nothing. Fix $p\geq8$. Put
\[
 L_0=[1,p],\qquad L_1=[3p,4p-2],\qquad L=L_0\cup L_1,
 \qquad N=|L|=2p-1.
\]
The high set is the following disjoint union:
\begin{align*}
 H_0&=[6p,8p-2],& H_1&=[8p,10p-2],& H_2&=\{10p\},\\
 H_3&=[11p-1,12p-1],& H_4&=[13p+1,14p-2],\\
 H_5&=[14p,15p-1],& H_6&=\{16p\},& H_7&=[17p-1,18p-1].
\end{align*}
Write $H=\bigcup_{i=0}^7H_i$ and $G=L\cup H$. The degree-two high
coefficient set is $\mathcal B=[6p,24p-1]\setminus H$, explicitly
\begin{align*}
 \mathcal B={}&\{8p-1,10p-1,14p-1\}\cup[10p+1,11p-2]\cup[12p,13p]\\
 &\cup[15p,16p-1]\cup[16p+1,17p-2]\cup[18p,24p-1].
\end{align*}
For later labels, $\mathrm C_0=\{8p-1\}$ and
$\mathrm C_2=[10p+1,11p-2]$. These are names for coefficient subsets,
not subquotients asserted to split off.

Define a low product with values in formal degree-two symbols by
\begin{equation}\label{eq:lowproduct}
 \mu(a,c)=
 \begin{cases}
  \mathsf A_{a+c},&a,c\in L_0,\ a+c>p,\\
  \mathsf B_{a+c},&a,c\text{ lie in different }L_i,\ a+c\geq4p-1,\\
  0,&\text{otherwise}.
 \end{cases}
\end{equation}
The total offset of the strand is
\[
 \tau=4p^2+6p-1=\sum_{a\in L}a+10p-2,
 \qquad \sum_{a\in L}a=4p^2-4p+1.
\]
Let $\mathcal S_p$ be free over $\Z$ on all labels $S[E;c]$ with
$E\subseteq G$, $|E|=2p-2$, $c\in L$, and $\sum E+c=\tau$.
Let $\mathcal K_p$ have the analogous basis $K[E;c]$ with $c\in H$.
All exterior sets are ordered increasingly. The target is free on labels
\[
 D[E;\mathsf A_d],\quad D[E;\mathsf B_d],\quad K[E;b],
 \qquad |E|=2p-3,\quad \sum E+d=\sum E+b=\tau,
\]
where $d\in[p+1,2p]$ for type $\mathsf A$,
$d\in[4p-1,5p-2]$ for type $\mathsf B$, and $b\in\mathcal B$.
Only the appropriate offset is used in each label. Denote the D and K
target summands by $\mathcal D'_p$ and $\mathcal K'_p$.

For $E=\{e_0<\cdots<e_{2p-3}\}$, define
\begin{align}
 M_pS[E;c]={}&
 \sum_{\substack{i:e_i\in L\\\mu(e_i,c)\ne0}}
 (-1)^iD[E\setminus\{e_i\};\mu(e_i,c)]\notag\\
 &+\sum_{\substack{i:e_i\in H\\e_i+c\in\mathcal B}}
 (-1)^iK[E\setminus\{e_i\};e_i+c],\label{eq:MS}\\
 M_pK[E;c]={}&\sum_{\substack{i:e_i+c\in\mathcal B}}
 (-1)^iK[E\setminus\{e_i\};e_i+c].\label{eq:MK}
\end{align}
Thus $M_p:\mathcal S_p\oplus\mathcal K_p\longrightarrow
\mathcal D'_p\oplus\mathcal K'_p$ is completely specified over the integers.
There is no implicit projection in these definitions. Our quotient is
$\coker_{\Z}M_p$, and a source identity always refers to this entire map.

\section{The four-row target and its earlier representative}\label{sec:eta}

Set
\[
 G_p(a,j;c)=K[(L\setminus\{a,3p,3p+j\})\cup\{6p\};10p+c].
\]
The tracked endpoint vector is
\begin{equation}\label{eq:eta}
 \begin{split}
 \eta_p=(-1)^p\bigl(&2G_p(p-3,2;p-3)-G_p(p-2,2;p-2)\\
                  &+2G_p(p-2,1;p-3)-2G_p(p-3,1;p-4)\bigr).
 \end{split}
\end{equation}
These are four distinct original rows. Their coefficient offsets are
$11p-3,11p-2,11p-3,11p-4$, all in $\mathrm C_2$.
Exactly one row coefficient is odd. This is a statement about a vector,
not a nonvanishing result in the cokernel.

For completeness, define the larger D representative without a choice of
quotient coordinates. Put
\[
 T(a,j)=S[(L\setminus\{a,3p,3p+j\})\cup\{6p,10p\};a+j-2]
\]
and set $s_p=\sum w_j(a)T(a,j)$, with the only nonzero weights
\begin{center}
\begin{tabular}{lll}
\toprule
$j$ & $a$ & $w_j(a)$\\
\midrule
2 & $1\leq a\leq p-4$ & $(-1)^{a+1}$\\
2 & $a=p-3$ & $(-1)^{p+1}$\\
1 & $a=p-2$ & $2(-1)^{p+1}$\\
1 & $a=p-3$ & $2(-1)^p$\\
\bottomrule
\end{tabular}
\end{center}
Let $b_p^A$ and $b_p^B$ be respectively the type-A and type-B parts of
$M_ps_p$, and write $b_p=b_p^A+b_p^B$. This defines them explicitly
through \eqref{eq:lowproduct}--\eqref{eq:MS}; it is the representative
recorded in EXP-052 and derived uniformly in EXP-056. Define also
\[
 q_p=K[(L\setminus\{3p,3p+2\})\cup\{6p\};10p].
\]

\begin{proposition}[[D] Original representative transfer]\label{prop:transfer}
For every $p\geq8$,
\begin{equation}\label{eq:transfer}
 M_p(s_p+q_p)=b_p+\eta_p.
\end{equation}
Consequently $[b_p]=-[\eta_p]$ in $\coker_{\Z}M_p$.
\end{proposition}
\begin{proof}
Every coefficient of $s_p$ lies in $[1,p-3]$. Its $6p$ face vanishes,
while its $10p$ face has negative sign and survives. Hence
$M_ps_p=b_p+\gamma_p$ with
$\gamma_p=-\sum w_j(a)G_p(a,j;a+j-2)$.
In $q_p$, the first-low faces $a=1,\ldots,p-2$ survive with sign
$(-1)^{a-1}$, and the other first-low faces vanish in $H_3$.
Every present second-low face vanishes in $H_4$, and the $6p$ face
vanishes at coefficient offset $16p$. Thus
\[
 M_pq_p=\sum_{a=1}^{p-2}(-1)^{a-1}G_p(a,2;a).
\]
Adding this expression to $\gamma_p$ cancels the interval
$1\leq a\leq p-4$ and leaves exactly \eqref{eq:eta}.
\end{proof}

\section{A complete integral potential basis in one sector}\label{sec:basis}

First fix the exterior high set $\{6p,8p-4\}$, not the high set used
later for the final annihilator. Index first-low variables as $p-r$,
$0\leq r\leq p-1$, and second-low variables as $3p+u$,
$0\leq u\leq p-2$. There are exactly two possible source families:
\begin{align*}
 A_u(r,s)&=S[(L\setminus\{p-r,p-s,3p+u\})\cup\{6p,8p-4\};p+2+u-r-s],\\
 B_r(u,v)&=S[(L\setminus\{p-r,3p+u,3p+v\})\cup\{6p,8p-4\};3p+2+u+v-r].
\end{align*}
Here $r<s$, $u+2\leq r+s\leq p+u+1$ in the first family. In the
second, $u<v$ and
\[
 \max(0,u+v-p+4)\leq r\leq\min(p-1,u+v+2).
\]
Indeed, the coefficient offset is the sum of the three omitted low
variables minus $4p-2$. Three omitted first-low variables make it negative;
three omitted second-low variables make it exceed $4p-2$. The remaining
two types give exactly these bounds. Thus no original source in this
fixed-high sector has been omitted.

\begin{theorem}[[D] Integral potential basis]\label{thm:basis}
Let $d_D$ be the D component of $M_p$ restricted to this sector. Choose
arbitrary integral potentials $f_u(r)$ vanishing for $r\leq u+1$.
For $r<s$ put $\alpha_u(r,s)=f_u(s)-f_u(r)$. For $u<v$ put
$R=u+v$, $F=f_u-f_v$, and
\[
 \beta_r(u,v)=
 \begin{cases}F(r)-F(R+3),&R\leq p-4,\\F(r),&R\geq p-3.
 \end{cases}
\]
The source with actual coefficients
\[
 (-1)^{p+r+s+u}\alpha_u(r,s)\quad\hbox{on }A_u(r,s),\qquad
 (-1)^{p+r+u+v}\beta_r(u,v)\quad\hbox{on }B_r(u,v)
\]
parametrizes $\ker_{\Z}d_D$ uniquely. Its unit potentials, indexed by
$0\leq u\leq p-3$, $u+2\leq r\leq p-1$, form an integral basis of
rank $\binom{p-1}{2}$. Each basis source has coefficient height one and
at most $3p-5$ terms; its full $M_p$ boundary has at most seven K rows
and no D rows.
\end{theorem}

\begin{proof}
We give the signed reconstruction, including truncation boundaries.
Complement an increasing low exterior $E$ inside $L$, and orient its
omitted variables in increasing first-endpoint order followed by increasing
second-endpoint order. For missing positions $d_1,\ldots,d_k$ in $L$,
the unreordered complement sign is
\[
 (-1)^{k(N-1)-\sum d_i-k(k-1)/2}.
\]
For $k=3$, reversing the omitted first-low variables gives exactly the
two source signs in the theorem. Complementing contraction gives wedge
insertion, with common factor $(-1)^{|E|-1}=-1$ here. Therefore all
nonzero D equations, up to a common unit sign on each row, are
\begin{align}
 \alpha_u(r,s)-\alpha_u(r,t)+\alpha_u(s,t)&=0,\label{eq:Aclosure}\\
 \alpha_u(r,s)-\alpha_v(r,s)+\beta_r(u,v)-\beta_s(u,v)&=0.\label{eq:Bclosure}
\end{align}
Equation \eqref{eq:Aclosure} has $r<s<t$ and
$u+2\leq r+s+t\leq p+u+1$. Equation \eqref{eq:Bclosure} has
$r<s$, $u<v$, and
$\max(0,R-p+4)\leq r+s\leq R+3$. Inadmissible source coefficients
are zero. These equations exhaust the A and B products in
\eqref{eq:lowproduct}.

For fixed $u$, read $f_u(r)=\alpha_u(0,r)$ when $r\geq u+2$ and
set the remaining values to zero. Every admissible pair $0<r<s$ occurs
in the A equation with endpoint zero. It forces
$\alpha_u(r,s)=f_u(s)-f_u(r)$. Pairs below the lower sum bound have
both potentials zero. Pairs above the upper bound cannot occur in an A
equation. Conversely these differences solve every A equation.

For a B equation, $r+s\leq R+3\leq p+u+1$ since $v\leq p-2$,
so upper truncation cannot spoil the potential formula. Lower truncation
has both relevant potentials zero. Equation \eqref{eq:Bclosure} becomes
$\beta_r-F(r)=\beta_s-F(s)$. If $R\leq p-4$, the beta interval is
$0,\ldots,R+2$. The equation on $\{0,R+3\}$ forces
$\beta_0=-F(R+3)$, and the equations on $\{0,r\}$ give the displayed
formula. If $R\geq p-3$, write $a=R-p+4\geq1$. Its beta interval
is $a,\ldots,p-1$, and $a\leq u+2$. Both potentials vanish below
$a$, so the equations with endpoint zero give $\beta_r=F(r)$.
The remaining equations hold by subtraction. Reconstruction is unique
and uses no division. The distinguished alpha coordinates therefore give
an integral inverse and an identity coordinate minor.

For one unit potential, alpha has at most $p-1$ terms. Beta is supported
only on second pairs containing its active second index. Usually such a
pair contributes at most one term. At most one pair satisfies
$R=r_0-3$ and contributes a constant interval of at most $p-1$ terms.
Thus beta has at most $2p-4$ terms, and all coefficients have absolute
value one.

Finally, every $6p$ face vanishes because
$6p+L_0\subset H_0$ and $6p+L_1\subset H_1$.
The negative $8p-4$ face survives only for first-low coefficient 3
or second-low coefficients $3p,3p+1,3p+2$. A unit potential gives at
most one first case, on $r+s=p+u-1$. In the second case there are at
most three ordinary partners, determined by $r_0-u_0$, $r_0-u_0-1$,
$r_0-u_0-2$, and at most three rows from the exceptional partner
$r_0-u_0-3$. This proves the full seven-row bound.
\end{proof}

\begin{remark}
This is a complete basis theorem for one D kernel. It is not a normal form
of $\coker M_p$, nor a claim that the whole connecting presentation has
seven rows. In particular, subsequent K relations may kill its images.
\end{remark}

\section{The shifted potential operator}\label{sec:shift}

For the annihilator, move the last high variable to $h=8p-2$.
Let $f_u(r)$ be integral and vanish for $r\leq u$. Define $P(f)$ by
the following two families, omitting zero terms and combining equal labels:
\begin{align}
 &S[(L\setminus\{p-r,p-s,3p+u\})\cup\{6p,h\};p+u-r-s],\notag\\[-3pt]
 &\hspace{18mm}\text{weight }(-1)^{p+r+s+u}(f_u(s)-f_u(r)),
 \quad r<s,\ u\leq r+s\leq p+u-1;\label{eq:Palpha}\\
 &S[(L\setminus\{p-r,3p+u,3p+v\})\cup\{6p,h\};3p+R-r],\notag\\[-3pt]
 &\hspace{18mm}\text{weight }(-1)^{p+r+u+v}\beta_r(u,v),
 \quad u<v,\ \max(0,R-p+2)\leq r\leq\min(p-1,R),\label{eq:Pbeta}
\end{align}
where $R=u+v$, $F=f_u-f_v$, and
\begin{equation}\label{eq:shiftbeta}
 \beta_r(u,v)=\begin{cases}F(r)-F(R+1),&R\leq p-2,\\F(r),&R\geq p-1.
 \end{cases}
\end{equation}
The coefficient offset is now the omitted-low sum minus $4p$, so all
labels have total offset $\tau$. The bounds put their coefficients in
$L_0$ and $L_1$, respectively. Moving the last high variable does not
change the exterior ordering sign. We assert a construction for these
potentials, not a complete basis theorem for the shifted sector.

\begin{lemma}[[D] Full shifted boundary]\label{lem:shift}
The D component of $M_pP(f)$ is zero. Every $6p$ face vanishes. The
only other surviving faces are the negative $h$ faces with first-low
coefficient 1 or second-low coefficient $3p$.
\end{lemma}
\begin{proof}
The A and B equations retain the signs \eqref{eq:Aclosure} and
\eqref{eq:Bclosure}, with bounds now
\[
 u\leq r+s+t\leq p+u-1,
 \qquad \max(0,R-p+2)\leq r+s\leq R+1.
\]
In an A equation, a pair below its admissibility bound has both potentials
zero, and no pair exceeds the upper bound. Its differences telescope.
In a B equation, $r+s\leq R+1\leq p+u-1$, so alpha again has no
upper truncation. Its equation is $\beta_r-F(r)=\beta_s-F(s)$.
For $R\leq p-2$, formula \eqref{eq:shiftbeta} extends as zero at
the absent endpoint $R+1$. No endpoint in the equation can exceed it.
For $R\geq p-1$, the lower beta endpoint is $a=R-p+2\leq u$;
both potentials vanish below $a$, so the second formula extends there
by zero. This proves every D equation.

As before, the $6p$ products lie in $H_0$ or $H_1$.
For $c\in L_0$, $h+c$ is in $\mathcal B$ only at $c=1$,
where it equals $8p-1$. The remaining values lie in $H_1$.
For $c=3p+t\in L_1$, $h+c=11p-2+t$ survives only at $t=0$;
the remaining values lie in $H_3$. The low exterior has even size
$2p-4$, so the last high face has negative sign. No other original
boundary row occurs.
\end{proof}

Use the signed K row
\begin{equation}\label{eq:erow}
 e(r;u,v)=(-1)^{p+r+u+v}
 K[(L\setminus\{p-r,3p+u,3p+v\})\cup\{6p\};11p-2+u+v-r],
\end{equation}
only for admissible indices and degree-two offsets. Put
$x_{uv}=e(u+v;u,v)$ when $u<v$ and $u+v\leq p-1$.
For orientation, let $y_{uv}=x_{uv}$ if $u<v$,
$y_{uv}=-x_{vu}$ if $u>v$, and $y_{uu}=0$.
In particular, $x_{01}=e(1;0,1)$ is different from $e(2;0,1)$.

The C0 boundary in Lemma~\ref{lem:shift} lies on
$r+s=p+u-1$. Its normalized coefficient is $-\alpha_u(r,s)$.
The C2 boundary is diagonal: for $R\leq p-2$ its coefficient on
$x_{uv}$ is $F(R+1)-F(R)$; at $R=p-1$ it is $-F(p-1)$.
There is no such diagonal row for $R>p-1$. This last endpoint cannot
be replaced by an unrestricted difference formula.

\section{Reflection intervals and the uniform annihilator}\label{sec:annihilator}

\begin{lemma}[[D] Signed interval triangle]\label{lem:triangle}
Let $i,a,b$ be distinct nonnegative second indices, with $a<b$ and
$i+a+b=p-2$. Let $T_i(a,b)$ have the single nonzero potential
\[
 f_i(r)=1\quad\text{for }i+a+1\leq r\leq i+b.
\]
Then the exact original boundary is
\begin{equation}\label{eq:triangle}
 M_pP(T_i(a,b))=y_{ia}-y_{ib}.
\end{equation}
\end{lemma}
\begin{proof}
The interval lies strictly inside $i<r<p-1$ and its endpoints sum to
$p+i-1$. It is therefore invariant under $r\mapsto p+i-1-r$.
Every possible C0 pair has this sum, so its two potential values agree,
including pairs outside the interval. All C0 rows vanish.
For C2 write $g(t)=f_i(i+t)$. Lemma~\ref{lem:shift} gives the
oriented edge coefficient $g(v+1)-g(v)$. The interval indicator jumps
by $+1$ at $v=a$ and by $-1$ at $v=b$, and nowhere else.
There is no top-endpoint contribution because the interval ends below
$p-1$. This proves \eqref{eq:triangle} in the full target, not just
after quotienting by K relations.
\end{proof}

For $j=1,2$, put $k=p-2-j$. Define the three-interval potential
\begin{equation}\label{eq:Fj}
 \begin{aligned}
 (F_j)_0(r)&=1&&\text{for }j+1\leq r\leq k,\\
 (F_j)_j(r)&=-1&&\text{for }j+1\leq r\leq p-2,\\
 (F_j)_k(r)&=-1&&\text{for }k+1\leq r\leq p-2,
 \end{aligned}
\end{equation}
and set all other values to zero. The indices $0,j,k$ are distinct
for every $p\geq8$.

\begin{proposition}[[D] Twice-edge source]\label{prop:twiceedge}
For $j=1,2$, $M_pP(F_j)=2x_{0j}$.
\end{proposition}
\begin{proof}
The definition is $F_j=T_0(j,k)-T_j(0,k)-T_k(0,j)$. Since $j<k$,
Lemma~\ref{lem:triangle} gives
\[
 M_pP(F_j)=(x_{0j}-x_{0k})-(-x_{0j}-x_{jk})-(-x_{0k}+x_{jk})
          =2x_{0j}.
\]
All three interval sources have zero C0 boundary individually.
The auxiliary edge satisfies $j+k=p-2$, within the diagonal domain.
\end{proof}

Two original K-source corrections are required to recover the particular
four-row target, rather than just the diagonal edges. Let $\delta_{0a}$
be the unit potential at $(u,r)=(0,a)$, and, for $a=2,3$, let
\begin{equation}\label{eq:Qa}
 Q_a=K[(L\setminus\{p-a,3p\})\cup\{6p\};8p-2-a]
\end{equation}
with \emph{positive} unit coefficient.

\begin{lemma}[[D] Exact short corrections]\label{lem:corrections}
Put $B=P(\delta_{03})-Q_3$ and $D=P(\delta_{02})+Q_2$. Then
\begin{align}
 M_pB&=e(3;0,1)+x_{02}+e(3;0,2),\label{eq:Bboundary}\\
 M_pD&=x_{01}+e(2;0,1).\label{eq:Dboundary}
\end{align}
\end{lemma}
\begin{proof}
Lemma~\ref{lem:shift} gives diagonal boundary
$x_{0,a-1}-x_{0a}$ for $P(\delta_{0a})$.
Its only C0 pair is $(a,p-1-a)$, with distinct increasing endpoints
for $a=2,3$ and $p\geq8$. The normalized alpha value is $-1$;
its source sign is $(-1)^{2p-1}=-1$ and the high-face sign is negative.
Thus the actual C0 coefficient is $-1$.

The only first-low face of $Q_a$ surviving in $\mathcal B$ removes
$a+1$, has coefficient offset $8p-1$, and sign $(-1)^a$.
Indeed the omitted first-low variable $p-a$ is larger than $a+1$,
so the latter has exterior position $a$. This is the same C0 row.
The surviving second-low faces remove $3p+v$, $1\leq v\leq a$.
Their coefficient offset is $11p-2-a+v$, and their exterior sign is
$(-1)^{p+v-2}$. In the signed convention \eqref{eq:erow}, their
coefficient on $e(a;0,v)$ is $(-1)^a$. Every other low product lies
in $H$. The high face has coefficient offset $14p-2-a$; the two
values $14p-4$ and $14p-5$ both lie in $H_4$ when $p\geq8$.
There is therefore no discarded high-free row or hidden filler.

Consequently the source $J_a=P(\delta_{0a})+(-1)^aQ_a$ cancels C0 exactly.
Its boundary is
\[
 M_pJ_a=x_{0,a-1}+\sum_{v=1}^{a-1}e(a;0,v).
\]
Specializing to $a=3$ and $a=2$ proves the lemma.
\end{proof}

\begin{theorem}[[D] Uniform original-map annihilation]\label{thm:main}
For every integer $p\geq8$, the explicitly defined integral source
\begin{equation}\label{eq:V}
 V_p=P(F_2+2F_1-4\delta_{03}-4\delta_{02})+4Q_3-4Q_2
\end{equation}
satisfies $M_pV_p=2\eta_p$. Hence both $[\eta_p]$ and $[b_p]$ have
order dividing two in the full integral cokernel.
\end{theorem}
\begin{proof}
Changing \eqref{eq:eta} from raw to signed row labels gives
\[
 \eta_p=-2e(3;0,2)-x_{02}-2e(2;0,1)-2e(3;0,1).
\]
Equations \eqref{eq:Bboundary}--\eqref{eq:Dboundary} therefore give
the exact target-vector equality
\[
 \eta_p=x_{02}+2x_{01}-2M_pB-2M_pD.
\]
By Proposition~\ref{prop:twiceedge}, twice the right-hand side is
\[
 M_p\bigl(P(F_2)+2P(F_1)-4B-4D\bigr)=M_pV_p.
\]
Every term is an admissible original source and every boundary row has
been included. Proposition~\ref{prop:transfer} also gives the explicit
identity $M_p(2s_p+2q_p-V_p)=2b_p$.
\end{proof}

\begin{corollary}[[D] Base change with two invertible]\label{cor:basechange}
Let $R$ be any commutative coefficient ring in which $2$ is a unit.
In the cokernel of $M_p\otimes_{\Z}R$, the images of $\eta_p$ and
$b_p$ vanish, with explicit sources $V_p/2$ and
$s_p+q_p-V_p/2$, respectively.
\end{corollary}
\begin{proof}
Tensor the two original source identities with $R$ and multiply by the
unit $2^{-1}$. This argument concerns only the two tracked classes.
\end{proof}

\section{A relative parity functional and exact order two}\label{sec:parity}

All coefficients in this section lie in $\mathbb F_2$. Write the full map as
\[
 M_p(s,t)=(d_Ds,\ d_K^Ss+d_K^Kt).
\]
A functional on the K target need not annihilate every S column individually.
It suffices to annihilate $d_K^K$ and $d_K^S(\ker d_D)$ while taking
value one on $\eta_p$. This relative formulation is essential: no functional
on a truncated matrix is substituted for a statement about the full map.

\subsection{The reflection template and its twelve-row specialization}

Let $T$ be any unordered triple of distinct indices in $[0,p-2]$ whose
sum is $p-2$. Define
\[
 z^T_{uv}=
 \begin{cases}
  1,&\{u,v,p-2-u-v\}\text{ is the multiset }T,\\
  0,&\text{otherwise}.
 \end{cases}
\]
Invalid indices and diagonals have value zero. Write $z=z^T$ temporarily.
It is symmetric and has the reflection identity
\begin{equation}\label{eq:zreflection}
 z_{uv}=z_{vu}=z_{u,p-2-u-v}.
\end{equation}
Assign the row $e(r;u,v)$ the value $z_{uv}$ when $r=u+v$ or $u+v+1$,
and zero otherwise. On a C0 row
\[
 K[(L\setminus\{p-r,p-s,3p+u\})\cup\{6p\};8p-1],
 \qquad r<s,\quad r+s=p+u-1,
\]
assign value
\begin{equation}\label{eq:dreflection}
 d_u(r)=z_{u,r-u}+z_{u,r-u-1}.
\end{equation}
Set every other original K row to zero. This defines $\lambda_T$.

The C0 definition does not depend on which endpoint is used. If
$n=p-1-u$ and $t=r-u$, then $s-u=n-t$, and
\eqref{eq:zreflection} exchanges the two summands of $d_u(r)$ and
$d_u(s)$. At a reflected fixed point $r=s$, the summands coincide and
$d_u(r)=0$. Thus no repeated exterior variable is introduced. Possible
cancellations in \eqref{eq:dreflection} are retained. More precisely, write
$T=\{i,j,k\}$ with $i<j<k$. For each $u\in T$, let $v<w$ be its other
two entries. The C0 pairs are
\[
 (r,s;u)=(u+v,u+w+1;u),\qquad (u+v+1,u+w;u).
\]
The second pair is omitted if $w=v+1$: it is a fixed point and its two
values cancel modulo two. All endpoints lie in $[0,p-1]$, and distinct
$u$ give distinct omitted second-low variables. There are always six
distinct C2 rows, so the exact support size is
\begin{equation}\label{eq:template-support}
 |\supp\lambda_T|=12-\mathbf1_{j=i+1}-\mathbf1_{k=j+1}.
\end{equation}
It ranges from ten to twelve; for example $T=\{1,2,3\}$ at $p=8$ has
ten rows. No numerical support count is assumed in this argument.

For the tracked class choose $T=\{0,2,k\}$, where $k=p-4$, and write
$\lambda_K=\lambda_T$. Its six C2 rows are
\[
 e(S;u,v),\ e(S+1;u,v),\qquad
 (u,v)=(0,2),(0,k),(2,k),\quad S=u+v.
\]
Its six C0 rows have omitted endpoints
\begin{equation}\label{eq:twelveC0}
 \begin{gathered}
 (r,s;u)=(2,p-3;0),(3,p-4;0),(2,p-1;2),\\
 (3,p-2;2),(p-4,p-1;k),(p-3,p-2;k).
 \end{gathered}
\end{equation}
All twelve rows are distinct and admissible for $p\geq8$. The C0
omitted-low sum is $4p+1$, which verifies its target grading; C2 grading
follows from \eqref{eq:erow}. The two coefficient types are distinct.
Modulo two, \eqref{eq:eta} is just $e(2;0,2)$, so
\begin{equation}\label{eq:eta-pairing}
 \lambda_K(\eta_p)=z_{02}=1.
\end{equation}

\subsection{All original K sources}

\begin{lemma}[[D] Complete K-source annihilation]\label{lem:Kparity}
Every reflection template $\lambda_T$ annihilates $d_K^K$.
\end{lemma}
\begin{proof}
A source with no exterior high variable cannot reach a supporting row.
With at least two exterior highs, a low face retains too many highs,
while a high face has coefficient at least $12p$, larger than every
supported coefficient. With one high different from $6p$, a low face
retains the wrong high and a high face has none. It remains only to
consider exterior high set $\{6p\}$.

Its low exterior omits two elements of $L$. If their sum is $m$, grading
forces coefficient $4p-2+m$. Two omitted first-low endpoints $r<s$
would give $6p-2-r-s<6p$, not an original K coefficient.
One omitted first-low and one second-low variable gives
\[
 X(r,u)=K[(L\setminus\{p-r,3p+u\})\cup\{6p\};8p-2+u-r].
\]
The only excluded high value in its possible coefficient range is $8p-1$,
so this column exists precisely when $r\geq u$ or $r\leq u-2$.
If $r\leq u-2$, neither supported C2 endpoint is possible and the
other C0 endpoint would exceed $p-1$. If $r\geq u$, its possible
C0 face has reflected endpoint $s=p+u-1-r$ and value $d_u(r)$.
At $s=r$ the face is absent but its assigned value is zero. Its C2
faces have second indices $r-u$ and $r-u-1$, with total value again
$d_u(r)$; invalid or already omitted indices have value zero.
The two contributions cancel, and the high face is invisible.

Finally, two omitted second-low indices $u<v$, of sum $S$, force
coefficient $10p-2+S$. This is in $H$ precisely for $S=2$ or
$S\geq p+1$. In the first case both supported first endpoints 2 and
3 are present and have the same value $z_{uv}$. In the second neither
$S$ nor $S+1$ is an available endpoint. Other faces are invisible.
This exhausts every original K source.
\end{proof}

\subsection{Reachability and complete D kernels}

An S source can reach the functional only by deleting a high variable
from an exterior high set $\{6p,h\}$. If $b$ is a supported target
coefficient and $c\in L$ the source coefficient, then $h=b-c\in H$.
The complete intersection calculation is
\begin{center}\small
\begin{tabular}{lll}
\toprule
$b$ & $c\in L_0$ & $c\in L_1$\\
\midrule
$8p-1$ & $[7p-1,8p-2]$ & none\\
$11p-2$ & $\{10p-2,10p\}$ & $[7p,8p-2]$\\
$11p-3$ & $\{10p-3,10p-2,10p\}$ & $[7p-1,8p-3]$\\
\bottomrule
\end{tabular}
\end{center}
Thus only
\begin{equation}\label{eq:visible-highs}
 h\in[7p-1,8p-2]\cup\{10p-3,10p-2,10p\}
\end{equation}
can matter. In particular $10p-3$ is present: first-low coefficient
$p$ reaches $11p-3$. The D differential preserves the complete
exterior-high set. Its kernel therefore splits into kernels of the
individual high-set sectors, and all sectors outside
\eqref{eq:visible-highs} have zero pairing term by term.

Consider first $h=8p-d$, with $2\leq d\leq p+1$. A source omits
three low variables of sum $m$, with coefficient $m-4p+d-2$.
Three omitted first-low variables make this at most $d-p-5<1$;
three omitted second-low variables make it at least $5p+d+1>4p-2$.
The two exhaustive source families and their coefficients are
\begin{align*}
 \alpha_u(r,s):&\quad p+d-2+u-r-s,
 &u+d-2&\leq r+s\leq p+u+d-3,\quad r<s;\\
 \beta_r(u,v):&\quad 3p+d-2+S-r,
 &\max(0,S+d-p)&\leq r\leq\min(p-1,S+d-2),\quad u<v.
\end{align*}
They have the same omitted-variable patterns as the alpha and beta
sources of Section~\ref{sec:basis}. Alpha coefficients can only lie
in $L_0$, beta coefficients only in $L_1$.

\begin{proposition}[[D] Complete generalized potential reconstruction]\label{prop:complete-d}
Every D cycle in this sector is uniquely parametrized over $\mathbb F_2$
by potentials satisfying
\[
 f_u(0)=0,\qquad f_u(r)=0\text{ for }r<u+d-2.
\]
Its coefficients are
\begin{align}
 \alpha_u(r,s)&=f_u(r)+f_u(s),\label{eq:fullalpha}\\
 \beta_r(u,v)&=
 \begin{cases}
 F(r)+F(S+d-1),&S\leq p-d,\\
 F(r),&S\geq p-d+1,
 \end{cases}\qquad F=f_u+f_v.\label{eq:fullbeta}
\end{align}
In particular, for $d=2$ the coordinates $f_u(u)$ are free for every
$u\geq1$. Only $f_0(0)$ is excluded by the endpoint-zero convention.
\end{proposition}
\begin{proof}
The complete A equations are
\[
 \alpha_u(r,s)+\alpha_u(r,t)+\alpha_u(s,t)=0,
 \quad u+d-2\leq r+s+t\leq p+u+d-3.
\]
Set $f_u(r)=\alpha_u(0,r)$ for $r>0$ with $r\geq u+d-2$,
and zero otherwise. The pair's upper bound is at least $p-1$.
Every other admissible pair occurs in its A equation with endpoint zero,
which forces \eqref{eq:fullalpha}. Conversely these potentials solve
all A equations: a pair below the lower bound has both potentials zero,
and a pair above the upper bound cannot occur in an admissible triple.

The complete B equations are
\[
 \alpha_u(r,s)+\alpha_v(r,s)+\beta_r(u,v)+\beta_s(u,v)=0,
 \quad\max(0,S+d-p)\leq r+s\leq S+d-1.
\]
Since $S+d-1\leq p+u+d-3$, both alpha terms can be replaced by
their potential sums, including lower truncations. The equation becomes
$\beta_r+F(r)=\beta_s+F(s)$.

If $S\leq p-d$, the beta interval is $0,\ldots,b$ with
$b=S+d-2\leq p-2$. The equation on $\{0,b+1\}$ forces
$\beta_0=F(b+1)$; the other endpoint-zero equations force the first
formula in \eqref{eq:fullbeta}. It extends as zero at $b+1$ and solves
every B equation. If $S\geq p-d+1$, put $a=S+d-p\geq1$.
The beta interval is $a,\ldots,p-1$, possibly empty, and
$a\leq u+d-2$. Both potentials vanish below $a$, so the endpoint-zero
equations force $\beta_r=F(r)$; this extends by zero below $a$.
If $a>p-1$, both potentials vanish everywhere and there is no beta
source. Thus all truncations, empty intervals, and converses are checked.
The reconstruction reads original coordinates, not a rank pattern.
\end{proof}

The additional free coordinates at $d=2$ cannot be suppressed here.
Section~\ref{sec:shift} used a restricted integral construction with
$f_u(r)=0$ for $r\leq u$; nonvanishing requires the complete kernel,
including its previously unused diagonal coordinates. The first-low endpoint
zero used in the reconstruction is available even when the smallest
second-low index in $T$ is positive. It places no restriction on the triangle.

\subsection{Pairing the complete kernels and the remaining sectors}

For a supported C0 face in a generalized potential cycle, the source
coefficient is $d-1$ and $r+s=p+u-1$. Reflection makes its value
$(f_u(r)+f_u(s))d_u(r)$. Fixed points contribute zero. Every nonzero
potential has $r\geq u$, so its reflected endpoint is also available.
Summing all distinct pairs gives
\begin{equation}\label{eq:parity-equal}
 \sum_{u,r}f_u(r)d_u(r)
   =\sum_{u<v}z_{uv}\bigl(F(u+v)+F(u+v+1)\bigr).
\end{equation}
On the support of a triple template, $1\leq S=u+v\leq p-2$,
so both potential indices in this sum are available.

For C2, only $r=S$ and $r=S+1$ can pair. Their source coefficients
are $3p+d-2$ and $3p+d-3$. If $d=2$, only $r=S$ is admissible;
\eqref{eq:fullbeta} gives $F(S)+F(S+1)$. If $3\leq d\leq p$,
both endpoints are admissible, and their common beta constant, when
present, cancels. If $d=p+1$, only $r=S+1$ is admissible. The sole
possible nonzero potential is $f_0(p-1)$, so $F(S)=0$ and the same
expression remains. Thus C2 pairing equals \eqref{eq:parity-equal}.
It cancels C0 over $\mathbb F_2$. Every other K face has functional
value zero; it need not itself vanish.

It remains to treat $h=10p-3,10p-2,10p$. Sources omitting two first-low
and one second-low variable cannot occur: their omitted sum is at most
$6p-3$, giving coefficient at most $-2$ even for $h=10p-3$.
Three omitted first-low variables also cannot occur. Sources omitting
three second-low variables, if admissible, have second-low coefficients
and do not contribute to A rows. Therefore the only source family with
first-low coefficient omits one first-low and two second-low variables.
For fixed $u<v$ and $S=u+v$, write its coefficients as $a_r$.
Its coefficient offsets are respectively
\[
 p+1+S-r,\qquad p+S-r,\qquad p-2+S-r.
\]
The lower endpoint bounds are $r\geq S+1$, $r\geq S$, and
$r\geq\max(0,S-2)$.

For each existing $r>0$, removal of the first-low variable $p$ gives
an A equation on endpoints $\{0,r\}$. Its only possible other
contributor is $a_0$. In the first two high sectors $a_0$ is absent,
so these equations force every $a_r$ to zero. For $h=10p$ the same
argument applies when $S\geq3$. For the remaining possible sums
$S=1,2$, use the A row on endpoints $\{S,S+1\}$. Its coefficient
offset is $2p-3-S>p$ and its only contributors are
$a_S+a_{S+1}$, which therefore vanishes. These are exactly the two
supported C2 contributions. This also covers triple templates containing
$\{0,1\}$; the selected twelve-row template itself has sums
$2,p-4,p-2$. There is no C0 contribution in these three high sectors.
Other families may contribute to B equations but cannot change the A
constraints used here.

We have proved the following relative statement on the complete map.
\begin{proposition}[[D] Relative reflection functional]\label{prop:relative-functional}
For every distinct triple template $T$ above,
\[
 \lambda_Td_K^K=0,\qquad
 \lambda_Td_K^S(s)=0\quad\text{whenever }d_Ds=0.
\]
This statement applies to every triangle, including those with positive
smallest entry; Section~\ref{sec:family} uses it to establish independence.
\end{proposition}

\begin{theorem}[[D] Exact order of the tracked class]\label{thm:exact-order}
For every $p\geq8$, $\eta_p$ is nonzero in the cokernel of
$M_p\otimes\mathbb F_2$. Its class in $\coker_{\Z}M_p$ has exact
order two. The same is true of the class of $b_p$.
\end{theorem}
\begin{proof}
Suppose $\eta_p=M_p(s,t)$ over $\mathbb F_2$. Its zero D component
implies $d_Ds=0$. Lemma~\ref{lem:Kparity} and
Proposition~\ref{prop:relative-functional}, for $T=\{0,2,p-4\}$,
would give
\[
 1=\lambda_K(\eta_p)=\lambda_K(d_K^Ss+d_K^Kt)=0,
\]
contradicting \eqref{eq:eta-pairing}. An integral source for $\eta_p$
would reduce to a source modulo two, so its integral class is also nonzero.
Theorem~\ref{thm:main} proves twice the class is zero. Its order is
therefore exactly two. Proposition~\ref{prop:transfer} transfers the
conclusion to $[b_p]=-[\eta_p]$.
\end{proof}

\section{The quadratic triangle family and its splitting}\label{sec:family}

Return to integral coefficients for source identities, and put $n=p-2$.
Let
\[
 \mathcal T_p=\{(i,j,k):0\leq i<j<k,\ i+j+k=n\},\qquad
 \mathcal C_p=\coker_{\Z}M_p.
\]
For $T=(i,j,k)$ select the original signed row $x_T=x_{ij}$ from
\eqref{eq:erow}. Its first endpoint $i+j$ is at most $p-2$, and its
coefficient offset is $11p-2$, so it is an admissible original target.
All subsequent relation statements take place in $\mathcal C_p$, not in
a projected or isolated subpresentation.

\begin{proposition}[[D] Full original triangle sources]\label{prop:all-triangle-sources}
For every $T=(i,j,k)\in\mathcal T_p$, define
\begin{equation}\label{eq:WT}
 F_T=T_i(j,k)-T_j(i,k)-T_k(i,j),\qquad W_T=P(F_T).
\end{equation}
Here each $T_u(a,b)$ is the integral interval potential of
Lemma~\ref{lem:triangle}, and $P$ retains exactly the source labels and
signs in \eqref{eq:Palpha}--\eqref{eq:shiftbeta}. Then
$M_pW_T=2x_T$ in the full original target.
\end{proposition}
\begin{proof}
For distinct $u,a,b$ with $a<b$ and $u+a+b=n$, the interval
$[u+a+1,u+b]$ is nonempty, lies strictly above $u$, and ends at
$n-a\leq p-2$. Its endpoints sum to $p+u-1$. Thus every interval
in \eqref{eq:WT} satisfies all hypotheses of Lemma~\ref{lem:triangle},
regardless of whether $i=0$. The complete D boundary cancels by the
signed potential equations; every $6p$ face vanishes; reflection kills
C0; and the two interval jumps give precisely its oriented C2 edges.
Consequently
\[
 \begin{split}
 M_pW_T
  &=(x_{ij}-x_{ik})-(-x_{ij}-x_{jk})-(-x_{ik}+x_{jk})\\
  &=2x_{ij}.
 \end{split}
\]
All coefficients here are integers. No unsupported boundary row has been
dropped, and no subsequent K relation is needed for the equality.
\end{proof}

\begin{theorem}[[D] Independent triangle classes and exact count]\label{thm:quadratic}
For every $p\geq8$, the relation lattice of the classes
$([x_T])_{T\in\mathcal T_p}$ in $\mathcal C_p$ is exactly
$2\Z^{\mathcal T_p}$. In particular there is an injection
\begin{equation}\label{eq:family-injection}
 \iota:(\Z/2)^{\mathcal T_p}\lhook\joinrel\longrightarrow\mathcal C_p,
 \qquad \mathbf e_T\longmapsto[x_T],
 \qquad |\mathcal T_p|=\left\lfloor\frac{(p-2)^2+3}{12}\right\rfloor.
\end{equation}
\end{theorem}
\begin{proof}
An unordered edge $\{a,b\}$ belongs to at most one triple of sum $n$:
its third entry must be $n-a-b$. Thus different triangles have disjoint
unordered edge sets. By the C2 definition of the relative functionals,
\begin{equation}\label{eq:diagonal-pairing}
 \lambda_T(x_U)=\delta_{TU}\quad\text{over }\mathbb F_2.
\end{equation}
Suppose that $\sum_U m_Ux_U=M_p(s,t)$ integrally. Reducing the full
source equation modulo two, its zero D component gives $d_Ds=0$.
Proposition~\ref{prop:relative-functional} therefore forces every
$\lambda_T$ to evaluate to zero on the right-hand K component.
Equation~\eqref{eq:diagonal-pairing} gives $m_T\equiv0\pmod2$ for
every $T$. Conversely, every even coefficient vector is a relation by
Proposition~\ref{prop:all-triangle-sources}. This proves the exact
relation lattice and the injection, not merely nonzero target vectors.

To count the triples, let $q(n)$ count strictly increasing nonnegative
triples of sum $n$. For $n\geq6$, triples with smallest entry at least
two correspond bijectively to those of sum $n-6$, by subtracting two
from every entry. Those with smallest entry zero number
$\lfloor(n-1)/2\rfloor$, and those with smallest entry one number
$\lfloor(n-2)/2\rfloor-1$. Hence
\begin{equation}\label{eq:count-recurrence}
 q(n)-q(n-6)=n-3\qquad(n\geq6).
\end{equation}
The initial counts for $n=0,\ldots,5$ are $0,0,0,1,1,2$.
The function $\lfloor(n^2+3)/12\rfloor$ has those same values and
the same recurrence, since the two unrounded arguments differ by the
integer $n-3$. Induction in each residue class modulo six proves the
formula for all $n\geq0$, hence for $n=p-2$.
\end{proof}

\begin{corollary}[[D] A direct summand of integral two-torsion]\label{cor:splitting}
The injection \eqref{eq:family-injection} splits as a homomorphism of
abelian groups. Thus, for some unspecified abelian group $\mathcal C'_p$,
\[
 \mathcal C_p\cong
 (\Z/2)^{\lfloor((p-2)^2+3)/12\rfloor}\oplus\mathcal C'_p.
\]
This is an existence statement for a retraction, not an explicit uniform
formula for its values on all D rows or a complete cokernel computation.
\end{corollary}
\begin{proof}
Work first over $\mathbb F_2$. For each $T$, define a functional on
the subspace $\operatorname{im}d_D$ of the full D target by
\[
 \mu_T(d_Ds)=\lambda_T(d_K^Ss).
\]
It is well-defined: two choices of $s$ differ by a complete D cycle,
whose connecting image is annihilated by
Proposition~\ref{prop:relative-functional}. Extend $\mu_T$ linearly
to the entire D target, by completing a basis of $\operatorname{im}d_D$.
The target is finite-dimensional for each fixed $p$, so such an extension
exists. The functional $(\mu_T,\lambda_T)$ annihilates the full map:
on an S source its two values are equal and cancel over $\mathbb F_2$,
and on a K source the value is zero by Lemma~\ref{lem:Kparity}.

Reduction modulo two followed by these extended functionals therefore
defines a group homomorphism
$\rho:\mathcal C_p\longrightarrow(\Z/2)^{\mathcal T_p}$.
Every selected target $x_U$ has zero D component, so
\eqref{eq:diagonal-pairing} gives $\rho\iota=\mathrm{id}$.
Thus $\mathcal C_p=\operatorname{im}\iota\oplus\ker\rho$.
The extension is noncanonical; no global D-row certificate is asserted.
\end{proof}

\begin{proposition}[[D] Exact transfer of the tracked class]\label{prop:eta-edge}
The original source
\begin{equation}\label{eq:eta-edge-source}
 C=P(F_1)-2B-2D
   =P(F_1-2\delta_{03}-2\delta_{02})+2Q_3-2Q_2
\end{equation}
satisfies $M_pC=\eta_p-x_{02}$. Hence
$[\eta_p]=[x_{02}]$ and $[b_p]=-[x_{02}]$ in $\mathcal C_p$,
where $x_{02}$ is chosen for the triangle $(0,2,p-4)$.
\end{proposition}
\begin{proof}
The target identity in the proof of Theorem~\ref{thm:main} is
$\eta_p=x_{02}+2x_{01}-2M_pB-2M_pD$. Apply
$M_pP(F_1)=2x_{01}$ and then Proposition~\ref{prop:transfer}.
This proves a quotient equality with an explicit source for its difference,
not equality of the original target vectors $\eta_p$ and $x_{02}$.
\end{proof}

In particular there are at least $3,4,5,7$ independent integral
two-torsion classes at $p=8,9,10,11$, respectively, and their number
is quadratically unbounded. If two is invertible in a coefficient ring,
each selected class vanishes with source $W_T/2$; this gives no vanishing
claim for the complementary quotient. Agreement of these lower bounds
with earlier finite isolated ranks does not identify that isolated object
with $\mathcal C_p$, nor does it establish an upper bound.

\section{Local obstruction and exact validation}\label{sec:validation}

\subsection{A failed locality prediction}

Before the interval construction, EXP-058 explored original source columns
incident to the four endpoint rows, retaining all their boundary rows.
Each row-column-row round gives the next radius. At $p=8$, the complete
declared radius-two span excludes $2\eta_8$ even over $\mathbb Q$ [MV].
The decisive certificates are integer row functionals:
\begin{center}
\begin{tabular}{rrrrr}
\toprule
Radius & Source columns & Boundary rows & Incidences & Dual support\\
\midrule
1 & 47 & 240 & 306 & 14\\
2 & 330 & 1803 & 2669 & 40\\
\bottomrule
\end{tabular}
\end{center}
For each selected matrix $M_{\mathrm{loc}}$, the stored functional satisfies
$\lambda M_{\mathrm{loc}}=0$ and $\lambda(2\eta_8)=4$.
The radius-two functional has 35 D rows and five K rows, all nonzero
coefficients equal to $1$ or $-1$. A coefficient-first reconstruction of
the complete neighborhoods, a separate original differential, and a dual
mutation check verify the certificates independently. The experiment stopped
at $p=8$; its planned $p=9,10$ cases were not run.

If a full source $z$ satisfies $M_8z=2\eta_8$, then some column outside
the selected radius-two set must have nonzero coefficient in $z$ and
nonzero pairing with $\lambda$. Otherwise applying $\lambda$ gives
$0=4$. This elementary escaping-column conclusion is compatible with
Theorem~\ref{thm:main}: the explicit source is not restricted to that
neighborhood. The local dual is not a functional on the full cokernel,
and its even target pairing is not a characteristic-two nonvanishing
certificate.

\subsection{Campaigns and adversarial controls}

For Theorem~\ref{thm:basis}, EXP-059 checks every unit potential at
$p=8,\ldots,16$, giving 525 chains, and four frozen endpoint pairs
at each $p=17,\ldots,100$, giving 336 more [MV]. All 861 full original
boundaries agree with the seven-row formula. The producer uses both
left-to-right and independently encoded right-to-left differentials.
A separate coefficient-first source reconstruction and bitset differential
audit all 861 cases; the latter also crosschecks the literal differential
on the 525 small cases. Source support ranges from 9 to 197 and boundary
support from 3 to 7 in this declared campaign. The integral completeness
claim, however, follows from the coordinate reconstruction proof.

For Theorem~\ref{thm:main}, EXP-060 checks the fixed parameters
\[
 8,9,\ldots,20,25,32,50,64,100.
\]
At all 18 values, two original differential implementations verify the
potential components, both short corrections, and $M_pV_p=2\eta_p$ [MV].
The independent auditor reconstructs all source labels and coefficients,
verifies the full boundary by a separately encoded bitset method, and
performs 25 literal differential crosschecks on five sources at each
$p=8,\ldots,12$. The complete boundary has four K rows and no D rows.
Observed source support is 110 at $p=8$ and 3054 at $p=100$, with
coefficient height five at each checked value. These are finite support
observations, not a separately proved uniform counting formula.

Wrong signs and single-coefficient changes are rejected. Two pre-declaration
paper errors are retained explicitly. Omitting the $\delta_{02}/Q_2$
correction gives $V_{\mathrm{omit}}=V_p+4D$ and
\[
 M_pV_{\mathrm{omit}}-2\eta_p=4(x_{01}+e(2;0,1))\ne0.
\]
The earliest proposal also had the opposite $F_1$ sign, giving
$V_{\mathrm{early}}=V_{\mathrm{omit}}-4P(F_1)$ and
\[
 M_pV_{\mathrm{early}}-2\eta_p=4(e(2;0,1)-x_{01})\ne0.
\]
The rows are distinct, so both failures hold uniformly. They were corrected
on paper before the final hypothesis was committed, not hidden after an
experimental failure.

Each producer and audit completed within its declared 60-second,
one-CPU-process and 1-GiB private-memory cap. There is no new HNF,
Smith-form or full ambient-basis computation in these campaigns. No old
labelled HNF source, including the original $p=11$ source holdout, is
an input. Checking the new formula at $p=11$ is not described as a new
untouched holdout. The dedicated suites contain 12 tests for the local
obstruction and 20 each for the potential basis and annihilator; tests
write only temporary outputs.

\subsection{Full-sector parity verification}

EXP-061 checks the twelve-row functional and 2,123 complete-sector unit
potential chains at $p=8,\ldots,16,25,32,50,64,100$ [MV]. Every basis
element is checked at $p\leq12$; larger parameters use frozen first,
middle and last basis indices, together with the newly free $d=2$
diagonal coordinate $(u,r)=(1,1)$. All chains have zero complete D
boundary and equal C0 and C2 pairings. The original K inverse-incidence
enumeration at $p=8,\ldots,12$ finds 13 incident K columns at each
parameter and annihilates every one. The reachable high-sector lists
include $10p-3$. The producer's 21 dedicated tests include the free
diagonal cases $(u,r)=(2,2)$ and $(4,4)$ at $p=8$, where the two
individual pairings are both one and cancel.

The independent audit does not use the producer's potential formulas.
It enumerates every original S source in every reachable high sector
at $p=8,\ldots,12$, constructs all original D rows over $\mathbb F_2$,
and certifies that the K pairing lies in the row span of that complete
D matrix. Equivalently, the pairing vanishes on its entire kernel.
The totals are 65 sectors, 23,695 original S sources, 108,261 D rows
and 231,986 incidences. The saved original-D-row dual certificates
contain 1,070 terms in total. Complete kernel dimensions agree with
the potential reconstruction, including the additional $d=2$ freedom.

All 20 independent adversarial controls pass. They include removal of the
odd eta-pairing row, mutation of a support index, an explicit diagnostic
for the otherwise omitted $h=10p-3$ sector, and a functional that passes
a proper local column subset but fails an added original column. The
producer and auditor each finish within their separately declared
120-second, one-process, 1-GiB caps. These finite checks challenge the
full-sector argument, not just a local span; the proof of
Theorem~\ref{thm:exact-order} remains the source classification and
relative-functional contradiction above.

\subsection{All-triangle identities and independent sector certificates}

EXP-062 verifies all 49 triangle sources at $p=8,\ldots,14$ and the
first, middle and last lexicographic triangle at each of
$p=16,20,25,32,50,64,100$, for 70 full integral identities [MV].
The producer also checks seven complete small-parameter pairing matrices,
93 count identities at $p=8,\ldots,100$, and the five original
$\eta_p-x_{02}$ transfers at $p=8,\ldots,12$. Both original
differential encodings agree on every signed source and transfer boundary.
The source selection was frozen before computation; the count theorem at
$p=100$ does not mean that every source at that parameter was run.

The independent auditor reconstructs all 70 signed sources, all five
transfers, the pairings, counts and mutation controls without importing
producer mathematics. For every one of the 27 triangles at
$p=8,\ldots,12$, it additionally enumerates all original S columns
in every reachable high sector and gives exact $\mathbb F_2$ row-span
certificates on the complete D matrix. The totals are
\begin{center}\small
\begin{tabular}{lr}
\toprule
Complete-audit quantity & Count\\
\midrule
Distinct $(p,h)$ sectors & 65\\
Distinct original S sources & 23,695\\
Triangle-sector certificates & 364\\
S-source instances across certificates & 151,319\\
D-row instances across certificates & 713,511\\
D incidences across certificates & 1,528,426\\
Labelled original D-row dual terms & 5,550\\
\bottomrule
\end{tabular}
\end{center}
Instance counts repeat a physical sector for different triangle functionals;
they are not counts of distinct matrices or columns. All original D rows
are retained in each certificate. This tests complete-kernel annihilation,
not merely the producer's constructed potential chains.

The ten-row adjacency case $T=(1,2,3)$ at $p=8$ is included. Removing
a functional row is detected by an explicit original K column. Changed
source signs or coefficients have nonzero full unit-column residuals;
permanent tests also multiply the actually altered sources. Duplicate
and mirrored edge choices give singular pairing matrices. All 39 combined
tests pass, comprising 24 producer and 15 independent-audit tests, and
the independent audit replays byte-for-byte to temporary outputs.

Observed source support ranges from 38 to 13,594, with coefficient height
two throughout this campaign. These are finite observations, whereas
\eqref{eq:template-support} proves the ten-to-twelve-row functional bound
uniformly. Producer and auditor each finish within their own declared
120-second, one-process, 1-GiB caps. No dense full ambient matrix, integral
HNF/SNF or old labelled $p=11$ HNF source is used. Finite certification
and the uniform proof remain separate layers of evidence.

\section{Remaining obligations}

The uniform twice-source, nonvanishing and independent-family obligations
are closed: Theorem~\ref{thm:quadratic} and Corollary~\ref{cor:splitting}
give the quadratically growing direct summand for every $p\geq8$.
The tracked endpoint class is explicitly identified with one selected
triangle class. What remains [C] is the complementary group: the complete
cokernel, a matching upper bound, its free rank and any additional torsion
are not determined. The functionals are not asserted to exhaust the full
dual space. In particular, the theorem does not identify the full
presentation with an earlier isolated component or relative-completion
quotient, and it does not prove a lower-strand recurrence.

The mathematical premise imported from the earlier family is the integral
coefficient presentation itself. Independent encodings attack implementation,
sign and boundary errors but do not rederive that family-to-presentation
identification. Within the explicit presentation of Section~\ref{sec:matrix},
the proofs above are self-contained. The new potential basis, signed
interval triangle and relative parity proof provide explicit sources and
detectors for further reductions. The parity argument exhausts the sectors
visible to each functional; it does not assert that those sectors exhaust
the whole map. The splitting is a proved existence result, but computing
a uniform global D-row retraction remains a separate constructive task.

\appendix
\section{Evidence map and source persistence}

All evidence is in the public \Repo\ repository, under the Huneke--Wiegand
experiment tree. The following links lead to the primary declarations,
proofs, verdicts, source code and exact artifacts. Mathematical claims are
owned by those primary records, not by a resume or summary page.

\begin{center}\small
\begin{tabular}{>{\raggedright\arraybackslash}p{0.43\textwidth}>{\raggedright\arraybackslash}p{0.46\textwidth}}
\toprule
Claim in this paper & Primary record\\
\midrule
Original presentation and product rules &
\Exp{EXP-036-factor-two-torsion-anatomy}{EXP-036},
\Exp{EXP-054-full-source-boundary}{EXP-054}\\[3pt]
Proposition~\ref{prop:transfer} and four-row target &
\Exp{EXP-056-uniform-low-source}{EXP-056},
\Exp{EXP-057-four-row-kernel-normal-form}{EXP-057}\\[3pt]
Local radius-two exclusion and duals &
\Exp{EXP-058-local-endpoint-source}{EXP-058}\\[3pt]
Theorem~\ref{thm:basis}, 861-chain audit &
\Exp{EXP-059-potential-connecting-map}{EXP-059}\\[3pt]
Lemmas~\ref{lem:shift}--\ref{lem:corrections}, Theorem~\ref{thm:main},
Corollary~\ref{cor:basechange} &
\Exp{EXP-060-uniform-endpoint-annihilator}{EXP-060}\\[3pt]
Complete relative parity proof and Theorem~\ref{thm:exact-order} &
\Exp{EXP-061-uniform-parity-functional}{EXP-061}\\[3pt]
Theorem~\ref{thm:quadratic}, Corollary~\ref{cor:splitting},
Proposition~\ref{prop:eta-edge} and complete triangle audits &
\Exp{EXP-062-triangle-torsion-family}{EXP-062}\\
\bottomrule
\end{tabular}
\end{center}

The full EXP-060 source archive retains every original exterior label and
integer coefficient for all 18 declared parameters. The EXP-062 archive
similarly retains all 70 tested triangle sources and five transfer sources.
Both deterministic gzip encodings have zero timestamp and empty embedded
filename. Their compact manifests store compressed and raw SHA-256 hashes,
byte counts, parameter counts, and per-source hashes. The auditors decompress
the archives and independently reconstruct all labels; no mathematical
information is discarded by compression.

\begin{center}\small
\begin{tabular}{lrrl}
\toprule
Artifact & Raw bytes & Stored bytes & SHA-256 prefix\\
\midrule
EXP-060 source archive & 4,508,227 & 80,250 & \texttt{85f8ea0800f0e8b6}\dots\\
EXP-060 result manifest & 462,616 & 462,616 & \texttt{f57fa704d045a8a8}\dots\\
EXP-060 independent audit & \multicolumn{2}{c}{exact JSON certificate} & \texttt{3bcf3cabe5b0330f}\dots\\
EXP-062 source archive & 17,253,771 & 252,812 & \texttt{8475ffb305c1567e}\dots\\
EXP-062 independent audit & \multicolumn{2}{c}{exact JSON certificate} & \texttt{854405294fab1d88}\dots\\
\bottomrule
\end{tabular}
\end{center}
Archive prefixes in this table identify the compressed bytes.
The EXP-060 raw digest starts with \texttt{d70c6ce4}\allowbreak\texttt{d7b331ef};
the EXP-062 raw digest starts with \texttt{af39d545}\allowbreak\texttt{d2f664de}.
Prefixes are navigation aids,
not sufficient integrity checks. Reproduction must verify the full 64-digit
hashes in the manifest and audit certificate. The accompanying permanent
tests also verify deterministic gzip replay, all stored source hashes,
failure checkpoint preservation, and the original $b_p$ source transfer
at $p=8$.

\begin{thebibliography}{9}
\bibitem{CAOSMain}
F.~Santiba\~nez-Leal,
\emph{Frobenius Minimality, Minimum-Layer Uniqueness, and an Infinite Family
of Huneke--Wiegand Counterexamples}, version 0.23, CAOS Research preprint,
2026. \href{https://doi.org/10.5281/zenodo.22181972}{doi:10.5281/zenodo.22181972}.

\bibitem{BrunsHerzog}
W.~Bruns and J.~Herzog,
\emph{Semigroup rings and simplicial complexes},
Journal of Pure and Applied Algebra \textbf{122} (1997), 185--208.
\href{https://www.home.uni-osnabrueck.de/wbruns/brunsw/pdf-article/SimpSemi.published.pdf}{Author-hosted published paper}.
\end{thebibliography}
\end{document}
